%=============================================================================
% Wave 4, Iteration 16: Materials Science Update --- EXOTIC PLATFORMS
% AGENT: MATERIALS SCIENCE (W4i16, Agent 13) | Model: Opus 4.6
%
% INHERITED FROM W4i15 MatSci:
%   - Si3N4 (proven) > TiN (new, superconducting) > crystalline Si (7K) >
%     Nb3Sn (Phase III)
%   - Cascaded MEMS: 10 Si3N4 membranes + CQNC = 500x margin at Q_psi=10^9
%   - Phononic-bandgap TLS suppression: 10^4x loss reduction
%   - Phononic Q=5e10 closes Q_mech gap
%   - 4H-SiC BLOCKED (zero built-in stress). 3C-SiC retains path.
%   - O-doping co-primary for Phase I
%   - Nb2O5 TLS origin confirmed
%   - Three-track roadmap: $300-530K / 18 months
%   - c_s^2 -> 0 theorem: DFD scalar sector does not propagate as wave
%   - Passive Superiority Theorem: only passive observations survive
%
% W4i16 MANDATE:
%   1. Graphene MEMS as psi-gradient sensors
%   2. Diamond NV centers for tidal psi-gradient measurement
%   3. Superfluid helium-3 B phase coupling
%   4. Metamaterials for psi-EM coupling enhancement
%   5. 2D materials beyond graphene (hBN, MoS2, WSe2)
%
% Date: 20 March 2026
%=============================================================================

\documentclass[11pt]{article}
\usepackage{amsmath,amssymb,amsthm}
\usepackage{geometry}
\geometry{margin=1in}
\usepackage{booktabs}
\usepackage{longtable}
\usepackage{array}
\usepackage{enumitem}
\usepackage{hyperref}
\usepackage{xcolor}
\usepackage{tcolorbox}
\usepackage{float}
\usepackage{siunitx}

\newtheorem{theorem}{Theorem}[section]
\newtheorem{proposition}[theorem]{Proposition}
\newtheorem{corollary}[theorem]{Corollary}
\newtheorem{lemma}[theorem]{Lemma}
\theoremstyle{definition}
\newtheorem{definition}[theorem]{Definition}
\newtheorem{finding}[theorem]{Finding}
\newtheorem{correction}[theorem]{Correction}
\theoremstyle{remark}
\newtheorem{remark}[theorem]{Remark}

\newcommand{\ee}[1]{\times 10^{#1}}
\newcommand{\Meff}{M_{\rm eff}}
\newcommand{\Qpsi}{Q_\psi}
\newcommand{\dpsi}{\delta\psi_{\rm EM}}
\newcommand{\psigrav}{\psi_{\rm grav}}
\newcommand{\psiEM}{\delta\psi_{\rm EM}}
\newcommand{\lambdaone}{|\lambda-1|}
\newcommand{\lambdabound}{1.56\ee{-9}}
\newcommand{\Tc}{T_{\rm c}}
\newcommand{\lambdaL}{\lambda_{\rm L}}
\newcommand{\FOMpassive}{\mathrm{FOM}_{\rm passive}}
\newcommand{\FOMMEMS}{\mathrm{FOM}_{\rm MEMS}}
\newcommand{\lam}{|\lambda-1|}
\newcommand{\Qmech}{Q_{\rm mech}}
\newcommand{\SNR}{\mathrm{SNR}}
\newcommand{\confirmed}[1]{\textcolor{blue}{\textbf{CONFIRMED:} #1}}
\newcommand{\corrected}[1]{\textcolor{red}{\textbf{CORRECTED:} #1}}
\newcommand{\caution}[1]{\textcolor{orange}{\textbf{CAUTION:} #1}}
\newcommand{\honest}[1]{\textcolor{violet}{\textbf{HONEST ASSESSMENT:} #1}}
\newcommand{\killed}[1]{\textcolor{red!70!black}{\textbf{KILLED:} #1}}
\newcommand{\verdict}[1]{\textcolor{green!50!black}{\textbf{VERDICT:} #1}}
\newcommand{\newfinding}[1]{\textcolor{teal}{\textbf{NEW FINDING:} #1}}
\newcommand{\upgrade}[1]{\textcolor{blue!70!black}{\textbf{UPGRADE:} #1}}

\title{Wave 4 Iteration 16: Materials Science Update\\
\large Exotic Platforms Assessment:\\
Graphene, Diamond NV, Superfluid $^3$He,\\
Metamaterials, and 2D Materials Beyond Graphene}
\author{DFD Research Programme --- Materials Science Agent 13, W4i16}
\date{20 March 2026}

\begin{document}
\maketitle
\tableofcontents
\newpage

%=============================================================================
\section{Executive Summary: Top 5 Findings}
\label{sec:top5}
%=============================================================================

\begin{tcolorbox}[colback=blue!5!white, colframe=blue!50!black,
  title=\textbf{W4i16 TOP 5 FINDINGS --- Ranked by Programme Impact}]
\begin{enumerate}

\item \textbf{FINDING 1 (CRITICAL KILL): Graphene drumhead MEMS are
KILLED for DFD $\psi$-gradient sensing despite extraordinary
theoretical $Q$. Two independent no-go arguments: (a)~effective mass
$\Meff \sim 10^{-16}$\,kg is $10^{6}$--$10^{7}\times$ below the DFD
requirement, and (b)~experimental $Q$ remains $\leq 2.5\ee{5}$ at 10\,K,
a factor $10^{3}$ below the $\Qpsi$ threshold.}

Graphene has the highest theoretical dissipation dilution factor of any
known material. The amplification coefficient for a graphene triangular
lattice phononic crystal (PnC) is $\sim 2833\times$ that of Si$_3$N$_4$
(Zheng et al., Nanomaterials 2022), owing to graphene's
atomic thickness (0.335\,nm vs $\sim$30\,nm for Si$_3$N$_4$). Multiplied
by graphene's intrinsic $Q \sim 200$ at room temperature, the theoretical
soft-clamped $Q$ reaches:
\begin{equation}
Q_{\rm graphene,\,PnC}^{\rm theory} \approx 200 \times D_Q^{\rm graphene}
\sim 10^8\text{--}10^9
\label{eq:graphene_theory}
\end{equation}
at room temperature, potentially exceeding $10^{10}$ at cryogenic
temperatures.

However, two fatal barriers prevent DFD application:

\textbf{(a) Effective mass catastrophe.} A monolayer graphene drumhead
of diameter $d = 30$\,$\mu$m has:
\begin{equation}
\Meff = \rho_{\rm 2D} \times \pi (d/2)^2 \times \alpha_{\rm mode}
\approx 7.7\ee{-7}\,\text{kg/m}^2 \times 7.1\ee{-10}\,\text{m}^2
\times 0.25 \approx 1.4\ee{-16}\,\text{kg}
\label{eq:graphene_mass}
\end{equation}
where $\rho_{\rm 2D} = 7.7\ee{-7}$\,kg/m$^2$ is the areal mass density
of monolayer graphene. DFD $\psi$-coupling scales
as $F_\psi \propto \Meff$, requiring $\Meff \sim 10^{-10}$--$10^{-9}$\,kg
for sub-SQMS sensitivity. The graphene mass deficit is
$\sim 10^{6}$--$10^{7}\times$. Even a 100-layer graphite drumhead only
reaches $\Meff \sim 10^{-14}$\,kg ($10^{4}$--$10^{5}\times$ deficit),
and multilayer graphite rapidly loses the dissipation dilution advantage.

\textbf{(b) Experimental $Q$ gap.} The highest demonstrated graphene
resonator $Q$ at cryogenic temperatures is $\sim 2.5\ee{5}$ (graphene/hBN
heterostructure at $< 10$\,K; Bachtold group, Nano Lett.\ 2017). Arrays
of bare graphene drumheads achieve only $Q \sim 100$--$2400$ at room
temperature (Bunch et al., Nano Lett.\ 2011). No graphene PnC membrane
has been fabricated to date --- the theoretical predictions of $Q \sim
10^{8}$--$10^{9}$ remain entirely unrealised.

\killed{Graphene MEMS for DFD. Both the $\Meff$ deficit ($10^{6}\times$)
and the $Q$ gap ($10^{3}\times$ experimental vs theory) are independently
fatal. The materials hierarchy is unchanged: Si$_3$N$_4$ $>$ TiN $>$
crystalline Si $>$ Nb$_3$Sn. Graphene offers no path to DFD-relevant
MEMS performance.}

\item \textbf{FINDING 2 (NEW ASSESSMENT): Diamond NV centres cannot
detect $\psi$-waves (killed by $c_s^2 \to 0$ theorem) but offer a
narrow \textit{calibration} role for SQMS Phase~I. Projected
$\psi$-detection sensitivity: $|\lambda - 1| \sim 10^{3}$
--- hopelessly above SQMS bound. NICHE as engineering tool only.}

The $c_s^2 \to 0$ theorem (W4i15) kills all resonant
$\psi$-wave sensors, including any NV-based oscillatory detection
scheme. The \textit{static} $\psi$-gradient from a laboratory EM source
is suppressed by $G/c^4 \sim 8.2\ee{-45}$\,m/J, yielding for an
SRF cavity ($u_{\rm EM} \sim 10$\,J/m$^3$, $r = 0.1$\,m):
\begin{equation}
\delta a_\psi \sim |\lambda - 1| \times \frac{u_{\rm EM}}{r^2}
\times \frac{G}{c^4}
\sim 1.56\ee{-9} \times 10^3 \times 8.2\ee{-45}
\approx 1.3\ee{-51}\,\text{m/s}^2
\label{eq:nv_signal}
\end{equation}
This is 49 orders of magnitude below even the best NV-MEMS
hybrid accelerometer sensitivity ($\sim 10^{-3}$\,m/s$^2/\sqrt{\text{Hz}}$).

State-of-the-art NV magnetic sensitivity stands at 2.6\,pT/$\sqrt{\text{Hz}}$
(Chatzidrosos et al., Phys.\ Rev.\ Appl.\ 2021), with room-temperature
improvements to 18\,pT/$\sqrt{\text{Hz}}$ reported in 2025. Deep-sea
operation has been demonstrated (NSR 2025).

\textbf{The one genuine value} of NV centres to DFD: as a
\textit{room-temperature magnetic field mapper} for SQMS Phase~I
cavity magnetic hygiene. An NV scanning probe ($\sim$nm spatial
resolution, 300\,K operation, $\sim$pT sensitivity) can map residual
fields inside the SRF cryostat to separate vortex-induced losses from
the DFD $Q$-ratio signal. Budget: \$5--10K.

\verdict{Diamond NV centres are a \textit{support tool} for SQMS Phase~I
(magnetic field mapping at \$5--10K). They have zero $\psi$-detection
capability (49-order gap). Not a detection instrument.}

\item \textbf{FINDING 3 (CRITICAL CORRECTION): Superfluid $^3$He-B
does \textit{NOT} break time-reversal symmetry. The B phase is
time-reversal \textit{invariant}. Furthermore, the $c_s^2 \to 0$ theorem
and $G/c^4$ suppression kill any DFD coupling in both $^3$He phases.
Confirms and extends W4i10 kill (34-order gap).}

The W4i16 mandate asked whether the broken TRS of the $^3$He-B phase
could enhance DFD coupling. This premise contains a factual error:

\begin{itemize}[nosep]
\item \textbf{B phase:} Time-reversal \textit{invariant} topological
  superfluid (class DIII, $\mathbb{Z}$ invariant). Isotropic
  Balian--Werthamer gap: $\Delta_{\alpha i} = \Delta_0 R_{\alpha i}
  (\hat{n}, \theta)$. Surface Majorana states protected by TRS.
  Magnetic field \textit{destroys} this protection (Silaev \& Volovik,
  2010).
\item \textbf{A phase:} Chiral $p$-wave ($p_x + ip_y$), \textit{breaks}
  TRS. Anisotropic gap. Stabilised in thin films $< 1$\,$\mu$m
  (Levitin et al., Science 2013).
\end{itemize}

Even correcting to the A phase, DFD coupling is negligible:
the $c_s^2 \to 0$ theorem kills resonant acoustic--$\psi$ coupling;
the $G/c^4$ suppression applies to all EM-sourced $\dpsi$; and the
QUEST-DMC bolometer threshold ($> 1$\,keV) is $10^{39}$ orders above
the DFD $\psi$-signal in a laboratory volume ($\sim 10^{-36}$\,eV).

The $^3$He continuous time crystal (Nature Comms.\ 2025, Aalto) couples
to a mechanical mode but at energy scales $\sim 10^{25}\times$ above the
DFD signal.

\killed{Superfluid $^3$He for DFD (both A and B phases). The TRS
premise for $^3$He-B was factually incorrect. Even with TRS-breaking
A phase, the $c_s^2 \to 0$ theorem + $G/c^4$ suppression + prior
W4i10 kill (34-order gap) are independently fatal.}

\item \textbf{FINDING 4 (ASSESSED --- NARROW UTILITY): Metamaterials
cannot enhance $\psi$-EM coupling (fundamental bottleneck is $G/c^4$,
not EM field amplitude), but a phononic metamaterial vibration shield
directly addresses the W3i9 Channel~B dominant systematic. Cost:
\$10--30K. Isolation: $> 60$\,dB at MEMS defect mode frequency.}

A metamaterial can increase local $u_{\rm EM}$ by concentrating the EM
field, but the total stored energy $U = u_{\rm EM} \times V_{\rm mode}$
is bounded by input power and EM $Q$. An SRF cavity already achieves
$Q_{\rm EM} \sim 10^{10}$ and $U \sim 10^{-2}$\,J. A metamaterial
concentrator with $V_{\rm meta} \sim 10^{-9}$\,m$^3$ reaches
$u_{\rm meta} \sim 10^7$\,J/m$^3$ but $U_{\rm meta} = 10^{-2}$\,J
--- \textbf{identical total stored energy, no net gain.}

Recent advances (Nature Comms.\ 2026: graded phononic metamaterials
on silicon wafers with $> 10^5$ unit cells; piezoelectric superlattice
phonon--polariton bandgaps, Frontiers in Physics 2025) demonstrate
that phononic bandgap engineering at the MEMS frequency scale
($\sim 1$\,MHz) is mature.

\newfinding{Phononic metamaterial vibration shield for MEMS. Bandgap
centred at MEMS defect mode ($\sim 1$\,MHz), bandwidth $\sim 100$\,kHz,
vibration isolation $> 60$\,dB. This directly addresses the Channel~B
dominant systematic (mechanical vibration, W3i9). Mount MEMS chip on
phononic metamaterial substrate. Compatible with all three tracks.
Cost \$10--30K. Should be added to Phase~II MEMS hardware specification
as a standard component.}

\item \textbf{FINDING 5 (COMPREHENSIVE KILL): All 2D materials beyond
graphene (hBN, MoS$_2$, WSe$_2$) are KILLED for DFD MEMS.
Best demonstrated $Q$ values are $10^3$--$10^4\times$ below the
$\Qpsi$ threshold, and all share the $\Meff$ catastrophe.}

\begin{center}
\begin{tabular}{l c c c c}
\toprule
\textbf{Material} & \textbf{Best $Q$ (RT)} & \textbf{Best $Q$ (cryo)}
& \textbf{$\Meff$ (typ.)} & \textbf{Gap to $\Qpsi$} \\
\midrule
Graphene & 2,400 & $2.5\ee{5}$ (10\,K) &
  $10^{-16}$\,kg & $10^3\times$ ($Q$) + $10^6\times$ ($M$) \\
hBN & $\sim 1{,}000$ & $\sim 10^4$ (est.) &
  $10^{-15}$\,kg & $10^4\times$ ($Q$) + $10^5\times$ ($M$) \\
MoS$_2$ & 3,315 & $\sim 10^4$ (est.) &
  $10^{-15}$\,kg & $10^4\times$ ($Q$) + $10^5\times$ ($M$) \\
WSe$_2$ & 1,200 & $4.7\ee{4}$ (4\,K) &
  $10^{-15}$\,kg & $5\ee{3}\times$ ($Q$) + $10^5\times$ ($M$) \\
\midrule
\textbf{Si$_3$N$_4$ PnC} & $5.1\ee{7}$ & $> 10^9$ (proj.) &
  $10^{-10}$\,kg & \textbf{NONE} \\
\textbf{TiN PnC} & --- & $8\ee{6}$ (2.2\,K) &
  $10^{-10}$\,kg & $\sim 25\times$ ($Q$ only) \\
\textbf{Cryst.\ Si} & --- & $> 10^{10}$ (7\,K) &
  $10^{-15}$\,kg (1D) & geometry gap only \\
\bottomrule
\end{tabular}
\end{center}

The WSe$_2$ result ($Q = 4.7\ee{4}$ at 4\,K; Morell et al., Nano
Lett.\ 2016) remains the cryogenic record for any TMD. The MoS$_2$
room-temperature record ($Q = 3315 \pm 115$; Nat.\ Microsyst.\
Nanoeng.\ 2024) comes from optimised 2\,$\mu$m diameter devices.
A 2025 advance (Nat.\ Microsyst.\ Nanoeng.) demonstrates functional
MoS$_2$ NEMS arrays with mass-transfer printing and independent
electronic tunability --- addressing scalability but not $Q$.

The root causes of low $Q$ in 2D TMDs are surface adsorbates and
edge defects: extreme surface-to-volume ratios make contamination
the dominant loss mechanism. No PnC soft-clamping has been demonstrated
in any TMD. hBN encapsulation improves $Q$ by $\sim 10\times$ but
does not close the $10^3$--$10^4$ gap to $\Qpsi$.

\killed{ALL 2D materials (hBN, MoS$_2$, WSe$_2$) for DFD MEMS. The
$Q$ gap ($10^3$--$10^4\times$) and $\Meff$ gap ($10^5$--$10^6\times$)
are independently fatal. This confirms and extends the W4i14 ``2D
materials KILLED'' assessment to all members of the family.}

\end{enumerate}
\end{tcolorbox}


%=============================================================================
\section{Detailed Analysis}
\label{sec:detailed}
%=============================================================================

\subsection{Graphene: Theoretical Promise vs Experimental Reality}

The dissipation dilution theory predicts that graphene's atomic-scale
thickness should yield the highest $Q$ amplification coefficient of any
material. Zheng et al.\ (2022) performed structural optimisation of
graphene triangular lattice PnC and found:
\begin{itemize}[nosep]
\item $Q$ amplification coefficient for graphene PnC: $\sim 5\ee{6}$
  (optimised 7$\times$13 cell structure)
\item Compared to $\sim 1800$ for equivalent Si$_3$N$_4$ PnC
\item Ratio: $\sim 2833\times$
\item Optimal structure is moderate-size (7$\times$13), not the largest
  --- dilution effect saturates
\end{itemize}

This theoretical framework assumes perfect crystallinity, perfect
clamping, zero surface contamination, and successful PnC patterning
without introducing defects. None of these have been achieved
experimentally. A 2024 attempt to fabricate PnC soft-supported
graphene resonators (Nanomaterials 14, 130) reported limited success.

\textbf{Fundamental mismatch with DFD requirements:} Even if
$Q \sim 10^9$ were achieved, the effective mass problem has no
solution. A graphene membrane large enough to achieve ng mass would
need area $\sim$1\,m$^2$ (monolayer) or $\sim$1\,cm$^2$ with
$\sim 10^4$ layers, at which point it is simply graphite and the
dissipation dilution advantage vanishes.


\subsection{Diamond NV Centres: Calibration Tool, Not Detector}

NV centres are primarily \textit{magnetic} sensors. For DFD detection,
one needs a transduction mechanism from $\nabla\psi$ to magnetic field
change. The most promising is a mechanical intermediary: a proof mass
on a cantilever whose displacement modulates the NV--magnet distance
(an NV-read MEMS accelerometer).

Best demonstrated NV-MEMS accelerometer sensitivity: $\sim
100$\,$\mu$g/$\sqrt{\text{Hz}}$ at room temperature.

For tidal $\psi$-gradients (astrophysical, Sun/Moon), the DFD
modification is $|\lambda - 1| \times g_{\rm tidal} \sim 1.56\ee{-9}
\times 10^{-6} \approx 1.56\ee{-15}$\,m/s$^2$. Superconducting
gravimeters reach $\sim 10^{-12}$\,m/s$^2/\sqrt{\text{Hz}}$. NV
is $\sim 10^9\times$ less sensitive again.

\textbf{Engineering value:} For SQMS Phase~I, mapping the residual
magnetic field inside the SRF cavity cryostat is critical for
separating vortex-induced losses from the DFD $Q$-ratio signal. An
NV gradiometer (two NV sensors, $\sim 1$\,cm separation) achieves
gradient sensitivity $\sim 100$\,pT/m/$\sqrt{\text{Hz}}$, sufficient
to map field non-uniformity at the nT level required for Phase~I
magnetic hygiene.


\subsection{Superfluid $^3$He: Comprehensive No-Go}

The QUEST-DMC experiment (J.\ Low Temp.\ Phys.\ 2026, Lancaster)
demonstrates $^3$He-B bolometry with vibrating nanowire resonators
(micron and sub-micron diameter) and SQUID readout at $\sim 100$\,$\mu$K.
Their energy threshold ($> 1$\,keV) is for dark matter detection.

For DFD, the energy deposited by the $\psi$-field in a laboratory
volume is:
\begin{equation}
E_\psi \sim |\lambda - 1| \times u_{\rm EM} \times V \times
\frac{G}{c^4} \sim 1.56\ee{-9} \times 10 \times 10^{-3} \times
8.2\ee{-45} \approx 1.3\ee{-55}\,\text{J} \approx 10^{-36}\,\text{eV}
\label{eq:he3_signal}
\end{equation}
This is $10^{39}$ orders below the QUEST-DMC threshold. The gap is
equivalent to detecting a single photon with a thermometer that requires
a supernova to register.


\subsection{Metamaterial Engineering: Phononic Shield Design}

The one actionable metamaterial application is a phononic bandgap
vibration shield:

\begin{itemize}[nosep]
\item \textbf{Purpose:} Suppress mechanical vibration coupling to the
  MEMS defect mode from cryostat vibrations, pulse-tube cooler, and
  acoustic noise.
\item \textbf{Design:} Graded phononic metamaterial on silicon wafer
  (scalable fabrication demonstrated at $> 10^5$ unit cells). Bandgap
  centred at MEMS defect mode frequency ($\sim 1$\,MHz), bandwidth
  $\sim 100$\,kHz.
\item \textbf{Predicted performance:} Vibration isolation $> 60$\,dB
  within bandgap. Converts the W3i9 Channel~B dominant systematic
  (mechanical vibration) from a noise floor limitation to a
  non-issue.
\item \textbf{Integration:} Mount MEMS chip directly on phononic
  metamaterial substrate. Compatible with all three tracks
  (Si$_3$N$_4$, TiN, crystalline Si).
\item \textbf{Cost:} \$10--30K for design + fabrication (standard
  silicon microfabrication).
\end{itemize}


\subsection{2D Materials: Final Comprehensive Survey}

The 2D materials landscape for NEMS resonators in 2024--2026 shows
continued progress in scalability and tunability but no breakthrough
in $Q$:

\begin{itemize}[nosep]
\item \textbf{MoS$_2$ arrays} (Nat.\ Microsyst.\ Nanoeng.\ 2025):
  Mass transfer printing enables large-scale arrays with independent
  gate tunability. $Q \sim 10^2$--$10^3$ at room temperature.
\item \textbf{MoS$_2$ optimisation} (Nat.\ Microsyst.\ Nanoeng.\ 2024):
  Tradeoff between device scale and surface nonidealities. Optimum at
  $d = 2$\,$\mu$m gives $Q = 3315 \pm 115$.
  $Q \cdot f$ product: $\leq 2\ee{10}$\,Hz (vs $> 10^{15}$ for
  crystalline Si at 7\,K).
\item \textbf{WSe$_2$} (Morell et al., Nano Lett.\ 2016): Monolayer
  at $Q = 4.7\ee{4}$ at 4.2\,K remains the cryogenic record for any
  TMD.
\item \textbf{hBN}: No dedicated cryogenic MEMS study. Used primarily
  as encapsulant for graphene heterostructures. Expected $Q \sim
  10^3$--$10^4$ at 4\,K.
\end{itemize}

All 2D materials share two DFD-incompatible properties:
(1)~$\Meff \sim 10^{-15}$\,kg ($10^{5}\times$ below threshold), and
(2)~$Q$ limited by surfaces (no PnC soft-clamping demonstrated in any
TMD, $10^3$--$10^4\times$ below $\Qpsi$).


%=============================================================================
\section{Impact on Materials Hierarchy and Roadmap}
\label{sec:impact}
%=============================================================================

W4i16 produces \textbf{zero changes} to the established materials
hierarchy and three-track roadmap:

\begin{tcolorbox}[colback=green!5!white, colframe=green!50!black,
  title=\textbf{Materials Hierarchy (UNCHANGED from W4i15)}]
\begin{enumerate}[nosep]
\item \textbf{Si$_3$N$_4$ PnC} (Phase~II, proven): $Q > 10^9$ at 4\,K.
  Highest confidence.
\item \textbf{TiN PnC} (Phase~II, new): $Q > 10^9$ projected. Best
  integration (native conductivity + superconducting).
\item \textbf{Crystalline Si PnC} (Phase~II+, 7\,K): Material ceiling
  $Q > 10^{10}$. Geometry gap only.
\item \textbf{Nb$_3$Sn} (Phase~III SRF): Passive FOM for cavity
  diagnostics.
\end{enumerate}
\textbf{KILLED (cumulative):} 4H-SiC, graphene, hBN, MoS$_2$, WSe$_2$,
all 2D materials, superfluid $^3$He, levitated nanospheres. \\
\textbf{NICHE:} Diamond NV (calibration tool, \$5--10K), metamaterials
(phononic vibration shield, \$10--30K).
\end{tcolorbox}

\textbf{Two additions to the roadmap:}
\begin{enumerate}
\item \textbf{Phononic metamaterial vibration shield}: Add to Phase~II
  MEMS hardware spec. Budget: \$10--30K. Timeline: design in 2026~Q4,
  fabricate with Track~3 Si$_3$N$_4$ run.
\item \textbf{NV scanning probe}: Add to SQMS Phase~I site
  characterisation. Budget: \$5--10K. Timeline: 2027~Q1.
\end{enumerate}

Total roadmap budget adjustment: +\$15--40K (negligible vs \$300--530K
three-track baseline).


%=============================================================================
\section{Corrections to Prior Iterations}
\label{sec:corrections}
%=============================================================================

\begin{correction}[W4i16-C1: $^3$He-B TRS]
\corrected{The W4i16 mandate stated that $^3$He-B has broken
time-reversal symmetry. This is incorrect. $^3$He-B is time-reversal
\textit{invariant} (class DIII topological superfluid). The A phase
breaks TRS. This error was in the mandate, not in any prior
iteration's findings. No prior MatSci update made claims about $^3$He.}
\end{correction}

No corrections to prior W4i15 findings. All five W4i15 findings remain
valid and unchanged.


%=============================================================================
\section{Summary and Forward Look}
\label{sec:summary}
%=============================================================================

W4i16 surveyed five exotic material platforms for DFD application:
\begin{itemize}[nosep]
\item \textbf{Graphene MEMS:} KILLED ($\Meff$ deficit $10^6\times$
  + $Q$ gap $10^3\times$)
\item \textbf{Diamond NV:} NICHE (calibration tool, \$5--10K)
\item \textbf{Superfluid $^3$He:} KILLED ($c_s^2 \to 0$ + $G/c^4$
  + prior 34-order kill)
\item \textbf{Metamaterials:} NICHE (phononic vibration shield,
  \$10--30K; addresses Channel~B systematic)
\item \textbf{2D TMDs (hBN, MoS$_2$, WSe$_2$):} KILLED ($Q$ gap
  $10^3$--$10^4\times$ + $\Meff$ gap $10^5\times$)
\end{itemize}

The materials science survey for DFD MEMS is \textbf{comprehensively
converged}. All reasonable exotic material candidates have been assessed
and rejected. The three-track roadmap (Si$_3$N$_4$ / TiN / crystalline
Si) is definitive. Future MatSci iterations should focus exclusively on
\textit{fabrication and integration engineering} within the established
tracks, not further materials screening.

\textbf{Recommended focus for W4i17 MatSci:}
\begin{enumerate}[nosep]
\item Phononic metamaterial shield design specification
\item TiN PnC fabrication process flow (leveraging superconducting
  qubit foundries)
\item Crystalline Si PnC mask design for EPFL/DTU collaboration
\item NV probe specification for SQMS Phase~I magnetic hygiene
\end{enumerate}

\end{document}
